\subsection{HiF8 Quantization}

\subsubsection{HiF8 Data Format}

HiF8 is a proprietary 8-bit floating-point format developed by Huawei.
In contrast to conventional FP8 variants (E4M3/E5M2), which fix the exponent
bit-width, HiF8 introduces a variable-width \emph{Dot} field that dynamically
partitions the available bits between the exponent and the mantissa, thereby
relaxing the intrinsic trade-off between dynamic range and representational
precision.
Its three core properties are summarized below.

\paragraph{Wide Dynamic Range.}
Conventional FP8-E4M3 fixes the exponent at 4~bits with a bias of 7, which yields
a limited exponent span and frequently induces overflow or underflow in the deep
attention layers of DiT-based generative models.
By unifying the normal and denormal encoding modes within a single coherent
scheme, HiF8 extends its effective exponent range to $[-22,\,15]$, covering 38
distinct exponent values and approaching the 40-exponent coverage of FP16
($[-24,\,15]$).
The resulting dynamic range enables HiF8 to quantize activations and gradients
that span several orders of magnitude without resorting to fine-grained
per-token or per-block rescaling, which materially reduces the engineering
complexity and memory overhead of low-precision training pipelines.

\paragraph{Tapered Precision.}
Conventional FP8 formats are constrained by a fundamental trade-off between
precision and dynamic range: E4M3 sacrifices dynamic range to obtain higher
mantissa precision, whereas E5M2 makes the opposite compromise, and neither
configuration simultaneously satisfies both requirements of neural-network
training.
HiF8 resolves this trade-off by progressively narrowing the mantissa bit-width
as the magnitude of a value increases: a 3-bit mantissa is preserved within the
central exponent range $[-3,\,3]$ and is gradually reduced to a 1-bit mantissa
at the extremes of the representable range.
This tapered precision profile is well matched to the approximately Gaussian
distribution of neural-network weights and activations, concentrating encoding
capacity in the high-density region of the numerical distribution while
relinquishing precision only at values that occur with low frequency.

\paragraph{Redundancy-Free Encoding.}
The exponent field of HiF8 employs Sign-Magnitude encoding with an implicit
leading 1-bit in the magnitude, which ensures that the representable ranges of
exponent fields of different widths are strictly non-overlapping and that the
entire $[-22,\,15]$ exponent space is covered without gaps or duplicates.
By contrast, conventional FP8 formats inherited from IEEE~754 contain redundant
code points, including dual representations of $\pm 0$ and multiple NaN encodings,
which waste code points within an already constrained 8-bit budget.
In HiF8, every code point is mapped to a unique numerical value, yielding full
encoding utilization and providing a stronger foundation for quantization accuracy
under tight bit-width constraints.

% -----------------------------------------------------------------------
\subsubsection{Quantization Scheme}

Mixed-precision training requires a \emph{scaling factor} that maps
high-precision tensors (BF16/FP32) into the representable range of HiF8.
Let $A_{\max} = \max|\mathbf{X}|$ denote the maximum absolute value within a
data block $\mathbf{X}$, and let $\mathrm{HiF8}_{\max}$ denote the largest
value representable in HiF8.
The scaling factor and the corresponding quantization operation are defined as

\begin{equation}
  \mathrm{Scale} = \frac{\mathrm{HiF8}_{\max}}{A_{\max}},
  \label{eq:hif8_scale}
\end{equation}

\begin{equation}
  \hat{\mathbf{X}} = \mathrm{cast\_to\_hif8}\!\left(\mathbf{X} \times \mathrm{Scale}\right).
  \label{eq:hif8_quant}
\end{equation}

\noindent\textbf{Current Scaling.}
We adopt a \emph{current scaling} strategy.
At each training iteration, the input tensor $\mathbf{X}$ is traversed to compute
the real-time value of $A_{\max}$, from which $\mathrm{Scale}$ is derived and
subsequently applied prior to the downstream computation (e.g., matrix multiplication).
This procedure guarantees that the scaling factor is consistent with the actual
data distribution at the current step and eliminates the overflow risk that would
otherwise arise from stale statistics.

\noindent\textbf{Per-Tensor Quantization Granularity.}
Owing to the wide dynamic range of HiF8, a single per-tensor scale is sufficient
for the entire tensor, and finer granularities such as the $128\!\times\!128$ tiles
typically required by FP8 formats are not necessary.
Per-tensor quantization substantially reduces the storage and memory-bandwidth
overhead associated with scale metadata, simplifies the control logic, and
contributes to higher end-to-end training throughput.

\noindent\textbf{Experimental Results.}
In training experiments on a DiT-based generative model (OSP-Next), the combination
of HiF8, current scaling, and per-tensor quantization, used in conjunction with the
Skiparse sparse-attention mechanism, yields loss curves and gradient statistics that
closely match those of the BF16 full-precision baseline.
These results indicate that HiF8 effectively balances numerical fidelity, training
stability, and system efficiency, and constitutes a practically viable 8-bit
mixed-precision solution for large-scale AIGC model training.
